\par\Needspace{9\baselineskip}\pdfbookmark[1]{Über die Summe der Reihe}{sixth-section-3}\begin{center}
\textit{DRITTER ABSCHNITT.}\\[6pt]
\textit{Über die Summe unserer Reihe, wenn das vierte Element $=1$ gesetzt wird, wobei zugleich einige andere}\\
\textit{transzendente Funktionen untersucht werden.}\\[6pt]
\end{center}
\SixthArticle{15}

Sind die Elemente $\alpha,\beta,\gamma$ sämtlich positive Größen, so sind alle Koeffizienten der Potenzen des vierten Elements $x$ positiv. Ist dagegen das eine oder andere dieser Elemente negativ, so haben wenigstens von irgendeiner Potenz $x^m$ an alle Koeffizienten dasselbe Vorzeichen, sofern $m$ größer als der absolute Wert des größten negativen Elements gewählt wird. Daraus ergibt sich ferner unmittelbar, dass die Summe der Reihe für
\SourceLeaf{148}
$x=1$ nur dann endlich sein kann, wenn die Koeffizienten wenigstens von einem gewissen Glied an gegen Null abnehmen, oder, um in der bei den Analysten üblichen Weise zu sprechen, wenn der Koeffizient des Gliedes $x^{\infty}$ $=0$ ist. Wir werden nun, und zwar mit voller Strenge zugunsten derjenigen, die den strengen Methoden der alten Geometer den Vorzug geben, zeigen,

\Needspace{14\baselineskip}
\textit{erstens}, dass die Koeffizienten (sofern die Reihe nicht abbricht) ins Unendliche wachsen, wenn $\alpha+\beta-\gamma-1$ eine positive Größe ist;

\textit{zweitens}, dass die Koeffizienten sich stetig einer endlichen Grenze nähern, wenn $\alpha+\beta-\gamma-1=0$ ist;

\textit{drittens}, dass die Koeffizienten gegen Null abnehmen, wenn $\alpha+\beta-\gamma-1$ eine negative Größe ist;

\textit{viertens}, dass die Summe unserer Reihe für $x=1$ trotz der Konvergenz im dritten Fall unendlich ist, wenn $\alpha+\beta-\gamma$ eine positive Größe oder $=0$ ist;

\textit{fünftens}, dass die Summe dagegen \textit{endlich} ist, wenn $\alpha+\beta-\gamma$ eine negative Größe ist.

